\section{Binding Flaws}\label{sec:definition}

A ZK application uses a proof the way a protocol uses a signature: as
permission to perform an action. The verifier certifies a \emph{statement}.
The consumer, however, performs an \emph{effect}. A binding flaw is a gap
between the two. This section defines the gap precisely and classifies the
ways it arises.

\subsection{Two ledgers}\label{sec:def-ledgers}

Consider a consumer function $f$ that, on input calldata $x$, sender $s$ and
consumer state $\sigma$, extracts a public-input vector $\vec{u} =
\mathsf{pi}(x,\sigma)$, calls $\mathsf{Verify}(\mathit{vk},\vec{u},\pi)$, and,
on acceptance, commits a state transition $\sigma' = \delta(x,s,\sigma)$.

\paragraph{Statement ledger}
The \emph{statement ledger} $\mathcal{S}(\vec{u})$ is the set of facts an
accepting proof establishes. It contains the public inputs $\vec{u}$ and every
value the circuit's relation $R$ uniquely determines from $\vec{u}$, such as a
recipient address hashed into a public commitment.

\paragraph{Effect ledger}
The \emph{effect ledger} $\mathcal{E}(f)$ lists every value that $\delta$ uses
to decide \emph{who} receives or is affected, \emph{how much} moves,
\emph{what} is recorded, \emph{which} state or policy applies, and \emph{how
often} the effect may occur. Recipients, amounts, fees, refund addresses,
nullifier keys, root identifiers, chain and contract identifiers, and
verification-key selectors all belong to the effect ledger.

\begin{figure}[t]
\centering
\begin{tikzpicture}[font=\scriptsize, box/.style={draw,rounded corners,minimum height=4.6cm,minimum width=2.9cm,align=center}, row/.style={draw,minimum width=2.5cm,minimum height=0.42cm}]
\node[box,label={[font=\small\bfseries]above:Statement ledger $\mathcal{S}$}] (S) at (0,0) {};
\node[box,label={[font=\small\bfseries]above:Effect ledger $\mathcal{E}$}] (E) at (4.6,0) {};
\foreach \i/\y/\t in {1/1.6/nullifier,2/0.9/Merkle root,3/0.2/amount,4/-0.5/{commit(recipient)}} \node[row] (s\i) at (0,\y) {\t};
\foreach \i/\y/\t in {1/1.6/{spent[key]},2/0.9/{root accepted?},3/0.2/{transfer amount},4/-0.5/{transfer to},5/-1.2/{fee, relayer},6/-1.9/{chain / contract}} \node[row] (e\i) at (4.6,\y) {\t};
\draw[-{Latex}] (s1) -- node[above]{bound} (e1);
\draw[-{Latex}] (s2) -- (e2);
\draw[-{Latex}] (s3) -- (e3);
\draw[-{Latex},dashed,red] (s4) -- node[above,red]{not checked} (e4);
\node[red,align=center] (att) at (1.2,-1.55) {attacker-chosen\\(no entry in $\mathcal{S}$)};
\draw[-{Latex},red] (att.east) -- (e5.west);
\draw[-{Latex},red] (att.east) -- (e6.west);
\end{tikzpicture}
\caption{The two ledgers of a ZK consumer. Binding holds when every
effect-ledger value is determined by the statement ledger and the consumer's
own state. Each red edge is a binding flaw.}
\label{fig:ledgers}
\end{figure}

\paragraph{Definition}
A consumer is \emph{bound} if for every accepting call, every value in
$\mathcal{E}(f)$ is a function of $\mathcal{S}(\vec{u})$ and the consumer
state $\sigma$ alone. A \emph{binding flaw} is a reachable violation, that is,
an execution in which three conditions hold:
\begin{description}
\item[C1] a proof is accepted, or the verifier accepts what it should reject;
\item[C2] the consumer commits an effect (a transfer, mint, state write, or
grant of access);
\item[C3] some value $e \in \mathcal{E}(f)$ that decides who, how much, what,
which state, or how often is chosen by the transaction sender and is not fixed
by $\mathcal{S}(\vec{u})$ and $\sigma$.
\end{description}
Because valid proofs are public once broadcast, any party can resubmit an
observed proof with different calldata. A flaw that lets a third party
redirect an effect by front-running is therefore a binding flaw, not an
ordering bug.
\TODO{T1}{state the definition as a short formal property over $(R,\mathsf{pi},\delta)$ and show that C1--C3 is exactly its negation}

\paragraph{Why neither layer suffices}
The definition relates two artifacts written in different languages. The
circuit determines $\mathcal{S}$. The consumer determines $\mathcal{E}$ and
$\mathsf{pi}$. A circuit analyzer can prove that $R$ is fully constrained, yet
it cannot know which consumer values the relation was meant to fix. A contract
analyzer can follow taint from calldata to a transfer, yet it cannot know that
the recipient was supposed to be a public input. Binding is a property of the
pair, and it holds only when the developer's intent, the circuit and the
consumer all agree.

\subsection{Taxonomy}\label{sec:def-taxonomy}

The flaws we observe fall into five patterns. Each mirrors a classic
signature-misuse failure (Table~\ref{tab:taxonomy}). The analogy is
predictive: for every pattern, the remedy is the one that fixes the
corresponding signature bug. The counts come from the measurement study in
Section~\ref{sec:study}.

\begin{table}[t]
\centering
\caption{Binding-flaw patterns, their signature-misuse analogues, and their
frequency among the 103 binding findings of Section~\ref{sec:study}.}
\label{tab:taxonomy}
\footnotesize
\setlength{\tabcolsep}{3pt}
\begin{tabular}{@{}p{2.0cm}p{1.75cm}p{2.45cm}r@{}}
\toprule
\textbf{Pattern} & \textbf{Signature analogue} & \textbf{Unbound value (example)} & \textbf{n} \\
\midrule
P1 Unbound effect & unsigned field & recipient, amount, fee, stored record & 23 \\
P2 Context & missing domain separation & chain, contract, policy, prior state & 23 \\
P3 Replay / alias & nonce reuse & nullifier key, $x$ vs.\ $x{+}r$ & 15 \\
P4 Verifier integration & weak or unpinned key & vk, program ID, pairing result & 17\,(+16) \\
P5 Interface mismatch & format confusion & order, width, encoding & 7 \\
Other & -- & -- & 2 \\
\bottomrule
\end{tabular}
% src: dataset/labels-v1/CHARACTERIZATION.md (pattern x severity); P4 split: LABELS-v1.csv proof_system_soundness (17 N, 16 Y)
\end{table}

\paragraph{P1: Unbound effect (unsigned field)}
The consumer uses a caller-supplied recipient, amount, fee, refund address or
stored value that is neither a public input nor hash-committed into one.
\begin{lstlisting}[language=Java]
function claim(Proof p, uint root, uint nf, address to) {
  require(verifier.verify(p, [root, nf]));   // `to` is not in the statement
  require(!spent[nf]); spent[nf] = true;
  token.transfer(to, AMOUNT);                 // anyone can copy p and redirect
}
\end{lstlisting}

\paragraph{P2: Context (missing domain separation)}
The proof does not commit to the chain, contract, configuration, policy or
\emph{prior state} in which the consumer uses it, so a proof made for one
context is accepted in another. Accepting a proof about a stale or
attacker-chosen state root falls here.

\paragraph{P3: Replay and alias (nonce reuse)}
The replay key is missing, is not derived from the statement (for example,
$\mathsf{keccak}(\pi)$ instead of a nullifier, which a re-randomized proof
evades), or aliases through the field: a contract that stores
\texttt{spent[nf]} on the raw 256-bit integer while the verifier reduces
inputs modulo $r$ accepts $nf$ and $nf+r$ as distinct.

\paragraph{P4: Verifier integration (weak or unpinned key)}
The deployed verifier, or the consumer's use of it, accepts statements it
should reject. Examples: unchecked pairing results, a verifier that returns
\texttt{true} regardless of the outcome, missing range checks on public inputs,
forgeable keys ($\gamma=\delta$), and a verification key or zkVM program ID
that the consumer does not pin. We count the 17 integration-level findings as
binding flaws. Another 16 findings are soundness bugs inside a proof system or
recursion verifier (transcript, FRI, or recursion-circuit constraints). They
have the same effect on the consumer, but their root cause is cryptographic,
so we report them as \emph{adjacent} and exclude them from the core class.
\TODO{W3}{author confirms the P4 split: 17 core plus 16 adjacent}

\paragraph{P5: Interface mismatch (format confusion)}
The consumer and the circuit disagree on the count, order, width or encoding
of public inputs, and the honest proof is still accepted while the consumer
reads a different value than the circuit proved. A mismatch that only makes
honest proofs fail is a liveness bug, not a binding flaw.

\paragraph{What is not a binding flaw}
Reentrancy and ordering issues where every effect value is fixed; denial of
service against one's own claim; trust in an administrator; and
circuit-internal under-constraints that do not change what the consumer is told
are all outside the class. One exception is that a proof consuming or
blocking \emph{another} user's state (e.g., spending a victim's nullifier) is
binding, because the affected party is not fixed by the statement.
